\section{Related Work}
\label{sec:related}

\paragraph{Entropy as a detector.} \textcite{shannon_ddos} detected network anomalies, such as denial-of-service attacks, from shifts in the entropy of traffic features like addresses and ports, without inspecting packet payloads or knowing the attack in advance. In LLMs, \textcite{shannon_hallucination} used semantic entropy to detect hallucinations. That work measures the uncertainty of one model answering one prompt.

\paragraph{Entropy in multi-agent systems.} \textcite{zhao2026doesmultiagentcollaborationhelp} measure entropy at the token, agent and round levels of collaborating agents and use it to predict when collaboration helps. In their setting communication is designed and known to the analyst, so entropy explains task performance rather than detecting coordination nobody intended.

\paragraph{Monitoring today.} AI control protocols place trusted monitors around untrusted models \parencite{greenblatt2024control}. These monitors read what the model produces, which stops working when messages are encoded or spread over resources nobody watches, as in the incident \parencite{openai2026incident, hf2026timeline}.

\paragraph{Our position.} We carry the network idea over to sandboxed agents. Agents take the place of network flows, their next-token distributions take the place of traffic features, and an unknown channel takes the place of an unknown attack. Unlike the work above, we use entropy to monitor for unintended coordination. The signal does not depend on the channel or the encoding, and it needs no extra model call, because the probabilities are produced anyway when each token is generated. We test it against a placebo that delivers useless foreign text and against cryptographic ground truth of whether communication happened.
